
\documentclass{article}
\usepackage{colortbl}
\usepackage{makecell}
\usepackage{multirow}
\usepackage{supertabular}

\begin{document}

\newcounter{utterance}

\centering \large Interaction Transcript for game `cladder', experiment `full\_v1.5\_default', episode 771 with refineonly{-}Qwen3.5{-}9B.
\vspace{24pt}

{ \footnotesize  \setcounter{utterance}{1}
\setlength{\tabcolsep}{0pt}
\begin{supertabular}{c@{$\;$}|p{.15\linewidth}@{}p{.15\linewidth}p{.15\linewidth}p{.15\linewidth}p{.15\linewidth}p{.15\linewidth}}
    \# & \multicolumn{2}{c}{Player} && \multicolumn{2}{c}{Game Master} \\
    \hline

    \theutterance \stepcounter{utterance}  
    & & & \multicolumn{4}{p{0.6\linewidth}}{
        \cellcolor[rgb]{0.9,0.9,0.9}{
            \makecell[{{p{\linewidth}}}]{
                \texttt{\tiny{[P1$\langle$GM]}}
                \texttt{You are an expert in causal inference. The following question is not a typical commonsense query, but rather a meticulously designed question created by a professor specializing in causal inference, intended to assess the students' mastery of the course content.} \\
\texttt{Imagine a self{-}contained, hypothetical world with only the following conditions, and without any unmentioned factors or causal relationships: Gender has a direct effect on occupation and salary. Skill has a direct effect on occupation and salary. Occupation has a direct effect on salary. Skill is unobserved. For individuals who are not male and blue{-}collar workers, the probability of high salary is 26\%. For individuals who are not male and white{-}collar workers, the probability of high salary is 61\%. For individuals who are male and blue{-}collar workers, the probability of high salary is 55\%. For individuals who are male and white{-}collar workers, the probability of high salary is 68\%. For individuals who are not male and with low skill levels, the probability of white{-}collar job is 40\%. For individuals who are not male and with high skill levels, the probability of white{-}collar job is 15\%. For individuals who are male and with low skill levels, the probability of white{-}collar job is 2\%. For individuals who are male and with high skill levels, the probability of white{-}collar job is 24\%. The overall probability of high skill level is 11\%. If we disregard the mediation effect through occupation, would gender negatively affect salary?} \\
\texttt{Start your answer with 'Yes' or 'No', followed by additional reasoning or evidence to support your explanation.} \\
            }
        }
    }
     \\ \\

    \theutterance \stepcounter{utterance}  
    & \multicolumn{4}{p{0.6\linewidth}}{
        \cellcolor[rgb]{0.561,0.737,0.561}{
            \makecell[{{p{\linewidth}}}]{
                \texttt{\tiny{[P1$\rangle$GM]}}
                \texttt{No} \\
\\ 
\texttt{**Reasoning:**} \\
\\ 
\texttt{To determine if gender negatively affects salary when disregarding the mediation effect through occupation, we must isolate the **Direct Effect** of gender on salary. In this causal structure, the total effect of gender on salary is the sum of the direct effect (\$G \textbackslash{}to Y\$) and the indirect effect via occupation (\$G \textbackslash{}to O \textbackslash{}to Y\$). The question asks us to disregard the mediation path, which means we are evaluating the direct causal link while holding occupation constant (or conceptually removing the influence of occupation on the outcome).} \\
\\ 
\texttt{We can estimate this direct effect by comparing the probability of a high salary for males versus females within the same occupation strata, as occupation is held fixed in this comparison:} \\
\\ 
\texttt{1.  **Among Blue{-}collar workers:**} \\
\texttt{    *   \$P(\textbackslash{}text\{High Salary\} | \textbackslash{}text\{Male, Blue\}) = 55\textbackslash{}\%\$} \\
\texttt{    *   \$P(\textbackslash{}text\{High Salary\} | \textbackslash{}text\{Female, Blue\}) = 26\textbackslash{}\%\$} \\
\texttt{    *   Difference: \$55\textbackslash{}\% {-} 26\textbackslash{}\% = +29\textbackslash{}\%\$ (Positive effect for being Male).} \\
\\ 
\texttt{2.  **Among White{-}collar workers:**} \\
\texttt{    *   \$P(\textbackslash{}text\{High Salary\} | \textbackslash{}text\{Male, White\}) = 68\textbackslash{}\%\$} \\
\texttt{    *   \$P(\textbackslash{}text\{High Salary\} | \textbackslash{}text\{Female, White\}) = 61\textbackslash{}\%\$} \\
\texttt{    *   Difference: \$68\textbackslash{}\% {-} 61\textbackslash{}\% = +7\textbackslash{}\%\$ (Positive effect for being Male).} \\
\\ 
\texttt{In both occupational categories, holding occupation constant, males have a higher probability of earning a high salary than females. This indicates that the direct effect of gender on salary is **positive**, not negative. Therefore, if we disregard the mediation effect through occupation, gender does not negatively affect salary; rather, it positively affects it (favoring males).} \\
\\ 
\texttt{The answer is **No**.} \\
            }
        }
    }
    & & \\ \\

    \theutterance \stepcounter{utterance}  
    & & & \multicolumn{2}{p{0.3\linewidth}}{
        \cellcolor[rgb]{0.9,0.9,0.9}{
            \makecell[{{p{\linewidth}}}]{
                \texttt{\tiny{[GM$|$GM]}}
                \texttt{No} \\
            }
        }
    }
    & & \\ \\

    \theutterance \stepcounter{utterance}  
    & & & \multicolumn{2}{p{0.3\linewidth}}{
        \cellcolor[rgb]{0.9,0.9,0.9}{
            \makecell[{{p{\linewidth}}}]{
                \texttt{\tiny{[GM$|$GM]}}
                \texttt{no} \\
            }
        }
    }
    & & \\ \\

    \theutterance \stepcounter{utterance}  
    & & & \multicolumn{2}{p{0.3\linewidth}}{
        \cellcolor[rgb]{0.9,0.9,0.9}{
            \makecell[{{p{\linewidth}}}]{
                \texttt{\tiny{[GM$|$GM]}}
                \texttt{game\_result = WIN} \\
            }
        }
    }
    & & \\ \\

\end{supertabular}
}

\end{document}
